% ============================================================
% Appendix Z.3 — FRFP Case Study (Justice / Governance):
% COMPAS Recidivism Risk Assessment
% ============================================================

\section{FRFP Case Study (Justice / Governance): COMPAS Recidivism Risk Assessment}
\label{app:frfp-compas}

This appendix analyses the COMPAS recidivism risk assessment tool as an
FRFP case study.
Rather than re-litigating the full fairness debate, we focus on how COMPAS
illustrates FRFP’s explicit/tacit split, boundary design, safe horizons,
and institutional non-automation in a high-stakes governance setting.

\subsection{Background}

COMPAS (Correctional Offender Management Profiling for Alternative
Sanctions) is a commercial risk assessment instrument originally developed
by Northpointe (now equivant) and used in parts of the U.S.\ criminal justice
system to assess the risk that a defendant will reoffend.
Jurisdictions have used COMPAS scores at various decision points, including
pre-trial release, sentencing, and parole.

In 2016, ProPublica published ``Machine Bias'', an investigation of COMPAS
scores and outcomes for over 7{,}000 defendants in Broward County, Florida.
They reported that:

\begin{itemize}
  \item overall accuracy (high vs.\ low risk vs.\ actual reoffending) was
        similar for Black and white defendants; but
  \item false positive rates were substantially higher for Black defendants
        (they were more often labeled high risk but did not reoffend), while
        false negative rates were higher for white defendants (more white
        defendants labeled low risk who did reoffend).
\end{itemize}

This sparked an extensive debate:

\begin{itemize}
  \item Northpointe and allied researchers responded that, under some
        standard statistical definitions (e.g.\ calibration conditional on
        risk score), COMPAS did not exhibit bias.
  \item Others argued that equal calibration and equal error rates across
        groups are mathematically incompatible in the presence of base-rate
        differences.
  \item Empirical studies showed that simple, transparent models or even
        human predictions based on a few factors can match COMPAS’s overall
        accuracy.
\end{itemize}

For FRFP, COMPAS is interesting not because of any single metric outcome,
but because it exemplifies:

\begin{itemize}
  \item explicit algorithms embedded within tacit judicial and institutional
        practices;
  \item ambiguous boundaries between ``advice'' and de facto decision;
  \item population-level effects over long horizons;
  \item attempts to automate governance via scores and dashboards.
\end{itemize}

\subsection{Explicit and Tacit Spaces in COMPAS Use}

\subsubsection*{Explicit space \(E_{\mathrm{justice}}\)}

COMPAS’s explicit domain includes:

\begin{itemize}
  \item input variables: answers to a questionnaire (137 items in some
        versions) on criminal history, demographics, and other factors;
  \item a proprietary scoring algorithm that maps inputs to risk scores and
        categories (e.g.\ low / medium / high risk);
  \item risk reports presented to judges or parole boards: numerical deciles,
        verbal risk levels, and explanatory text;
  \item policy documents and bench books describing how scores may be used.
\end{itemize}

All of this lives in explicit space: the algorithm is an explicit mapping;
scores are explicit artefacts; policy documents are explicit guidance.

\subsubsection*{Tacit space \(T_{\mathrm{justice}}\)}

Tacit space comprises:

\begin{itemize}
  \item judges’ and parole officers’ understanding of crime, risk, and
        rehabilitation;
  \item legal norms and intuitions about proportionality, due process, and
        equality;
  \item practical knowledge of local conditions (policing intensity,
        community context, available programmes);
  \item institutional culture around sentencing and pre-trial decisions.
\end{itemize}

These tacit states shape how explicit artefacts (COMPAS scores, case files,
arguments) are interpreted and integrated into decisions about liberty and
punishment.

\subsection{Boundary Design: Scores as Inputs vs.\ De Facto Decisions}

\subsubsection*{Formal role: advisory inputs}

Officially, risk assessment tools like COMPAS are often framed as providing
\emph{advisory} information:

\begin{itemize}
  \item judges retain legal authority and responsibility for sentencing or
        release decisions;
  \item COMPAS scores are one factor among many in a pre-sentence report;
  \item guidelines and court decisions (e.g.\ \emph{State v.\ Loomis} in
        Wisconsin) emphasise that COMPAS cannot be the sole basis for a
        decision and that its limitations must be acknowledged.
\end{itemize}

In FRFP terms, the intended boundary is:

\begin{itemize}
  \item AI / algorithm outputs (scores) remain in explicit space;
  \item human decision-makers in tacit space endorse or reject specific
        sentencing or release decisions, having weighed scores alongside other
        evidence.
\end{itemize}

\subsubsection*{Practical drift: the ``siren song of objectivity''}

In practice, several factors push toward a different boundary:

\begin{itemize}
  \item scores are presented as quantitative, objective indicators;
  \item judges face time pressure and may have limited statistical training;
  \item there is social and institutional pressure to appear ``data-driven''
        and to follow recommended tools;
  \item defendants and lawyers often have limited ability to challenge or
        interpret scores, given COMPAS’s proprietary nature.
\end{itemize}

The result is a drift toward treating COMPAS scores as:

\begin{itemize}
  \item strong \emph{priors} on risk that tacit judgment rarely overrides;
  \item anchors that bias subsequent interpretive processes.
\end{itemize}

From an FRFP perspective, this is boundary erosion:
explicit artefacts (scores) become de facto bearers of correctness, even
though legal doctrine maintains that humans remain responsible.

\subsection{Protocol-Level Issues in FRFP Terms}

\subsubsection*{Z.3.1 Explicit-Space Fairness and Its Limits}

ProPublica’s analysis and subsequent work exposed that COMPAS’s explicit
mapping from features to risk scores embodies trade-offs between accuracy
and fairness:

\begin{itemize}
  \item the tool can be approximately calibrated across groups (similar risk
        scores correspond to similar empirical recidivism rates);
  \item or it can equalise false positive and false negative rates across
        groups;
  \item but it cannot achieve both simultaneously when base recidivism rates
        differ between groups, as they do in the Broward dataset.
\end{itemize}

This is now a standard result in algorithmic fairness:
calibration, equal false positive rates, and equal false negative rates
cannot, in general, be satisfied together.

In FRFP language:

\begin{itemize}
  \item many fairness properties of interest are properties of explicit
        mappings in \(E_{\mathrm{justice}}\) and their empirical performance;
  \item explicit semantics are constrained by statistical structure in the
        environment (base rates), leading to impossibility trade-offs;
  \item no choice of metric can fully capture tacit normative judgments about
        what counts as just or unjust treatment across groups.
\end{itemize}

Thus, COMPAS illustrates that explicit-only fairness optimisation has
built-in limits: the algorithm can be tuned for one metric, but cannot
resolve deeper normative conflicts.

\subsubsection*{Z.3.2 Boundary Confusion and Responsibility}

The \emph{Loomis} decision in Wisconsin, which upheld the use of COMPAS
scores but warned about their limitations, is an example of FRFP-relevant
boundary tension.

From an FRFP viewpoint:

\begin{itemize}
  \item courts attempt to keep correctness in tacit space
        (judicial judgment, due process, constitutional norms);
  \item yet they allow opaque explicit artefacts (scores from proprietary
        algorithms) to enter the decision pipeline;
  \item defendants cannot meaningfully contest or interpret the explicit
        mapping, undermining the transparency usually required for legal
        correctness.
\end{itemize}

This creates a misalignment:

\begin{itemize}
  \item tacit responsibility is formally with judges;
  \item but some of the effective decision-making power has migrated into
        explicit systems that are not fully intelligible to any actor in
        \(T_{\mathrm{justice}}\).
\end{itemize}

FRFP would flag this as a violation of clear boundary design: humans are
supposed to be endorsers of explicit artefacts, but here the artefact is
partly inscrutable.

\subsubsection*{Z.3.3 Safe Horizons and Population Effects}

COMPAS is used repeatedly across thousands of cases over many years:

\begin{itemize}
  \item each use updates tacit beliefs of judges, lawyers, and communities
        about who is ``high risk'';
  \item scores may shape plea negotiations, sentencing norms, and incarceration
        patterns, feeding back into future data.
\end{itemize}

FRFP emphasises that in such long-horizon, population-level settings:

\begin{itemize}
  \item tacit grounding (e.g.\ nuanced understanding of defendants and local
        conditions) can erode as actors defer to the tool;
  \item safe horizons for tool usage should be explicitly defined and
        periodically reset via audits and moratoria;
  \item otherwise, structural biases can become entrenched.
\end{itemize}

COMPAS deployments generally did not incorporate such safe-horizon designs:

\begin{itemize}
  \item there were no built-in limits on how many cases could be decided with
        COMPAS involvement before formal reassessment of its effects;
  \item debates about bias and fairness arose only after sustained external
        pressure (e.g.\ ProPublica’s investigation, academic critiques).
\end{itemize}

From an FRFP perspective, this is a predictable pattern for explicit tools
that operate without structural safe horizons in governance contexts.

\subsubsection*{Z.3.4 Institutional Operators and Explicit-Only Governance}

At the institutional level, risk assessment is often framed as:

\begin{itemize}
  \item a means to improve consistency and transparency of decisions;
  \item a way to manage caseloads and resources more ``objectively'';
\end{itemize}

However, FRFP’s institutional semantics would ask:

\begin{itemize}
  \item which operators map the population’s tacit states (judges, lawyers,
        policymakers, communities) into new states as COMPAS is adopted?
  \item are there explicit governance bodies that:
        \begin{itemize}
          \item monitor fairness trade-offs over time;
          \item adjust boundaries and horizons;
          \item decide when to suspend or restrict the tool?
        \end{itemize}
\end{itemize}

The COMPAS story suggests that much governance has been:

\begin{itemize}
  \item reactive (responding to criticism and litigation);
  \item fragmented across jurisdictions;
  \item heavily reliant on explicit metrics and vendor documentation, rather
        than on a coherent FRFP-style tacit governance layer.
\end{itemize}

This fits FRFP’s warning about explicit-only governance in complex
institutional systems.

\subsection{Counterfactual: An FRFP-Aligned Risk Assessment Regime}

Rather than advocating for or against algorithmic risk assessment in
principle, FRFP emphasises structural conditions for any such system.

\subsubsection*{Z.3* Explicit/Tacit Role Clarification}

\begin{itemize}
  \item Explicit tools (including COMPAS-like models) are treated as providers
        of \emph{signals} and structured summaries, not as direct inputs into
        sentencing ranges or binary decisions.
  \item Tacit legal and ethical judgment (judges, parole boards, oversight
        bodies) retains responsibility for interpreting these signals in
        context.
\end{itemize}

\subsubsection*{Z.3* Boundary Design and Transparency}

\begin{itemize}
  \item Boundary events are clearly defined: each sentencing or release
        decision is an endorsement in \(T_{\mathrm{justice}}\), not in
        \(E_{\mathrm{justice}}\).
  \item For any decision in which an algorithmic score materially influenced
        the outcome, the record includes:
        \begin{itemize}
          \item the score and category;
          \item a human-readable explanation of its role;
          \item space for judicial justification, including reasons for
                overriding or discounting the score.
        \end{itemize}
  \item Defendants have structured avenues to challenge both the use of the
        tool and its score in their case.
\end{itemize}

\subsubsection*{Z.3* Safe Horizons and Audits}

\begin{itemize}
  \item Use of risk scores is subject to explicit safe horizons:
        \begin{itemize}
          \item after a certain number of decisions or time period, a full
                audit of outcomes and fairness metrics is mandatory;
          \item algorithm use can be paused or restricted pending audit.
        \end{itemize}
  \item Audits examine:
        \begin{itemize}
          \item group-wise error rates and calibration;
          \item downstream impacts on incarceration, bail, and recidivism;
          \item interaction with policing and charging patterns.
        \end{itemize}
  \item Trade-offs between fairness metrics are surfaced for public and
        legislative deliberation, not buried in model design.
\end{itemize}

\subsubsection*{Z.3* Institutional Governance}

\begin{itemize}
  \item A dedicated risk-assessment governance body (judges, defenders,
        prosecutors, community representatives, statisticians) oversees
        selection, deployment, and revision of tools.
  \item This body operates as an FRFP institutional operator:
        \begin{itemize}
          \item it defines admissible pipelines and boundary rules;
          \item it sets safe horizons and audit protocols;
          \item it retains the authority to withdraw tools that fail to satisfy
                structural or normative conditions.
        \end{itemize}
\end{itemize}

Under such a regime, COMPAS-like tools remain in explicit space; correctness
over citizens’ liberty remains a tacit, human and institutional function.

\subsection{Lessons for FRFP and Algorithmic Governance}

The COMPAS case reinforces several of FRFP’s core claims:

\begin{enumerate}
  \item \textbf{Fairness trade-offs are explicit and structural.}
        COMPAS shows that explicit mappings from features to scores are
        constrained by base rates and cannot satisfy all fairness desiderata
        simultaneously.
        This is not a bug of one model but a structural fact.

  \item \textbf{Boundaries must be designed, not assumed.}
        Declaring a tool ``advisory'' does not prevent its scores from
        shaping decisions as if they were authoritative.
        FRFP requires explicit, auditable boundary mechanisms.

  \item \textbf{Safe horizons are essential for population-level use.}
        Long-term deployment of risk tools without periodic re-grounding in
        tacit legal and ethical judgment risks entrenching bias and eroding
        human responsibility.

  \item \textbf{Explicit-only governance is inadequate.}
        Oversight based solely on vendor documentation and retrospective
        metrics cannot, in general, recover tacit notions of justice and due
        process; institutional FRFP operators are needed.

  \item \textbf{Formal architectures clarify debates.}
        The COMPAS controversy is often framed as ``is the tool racist or
        not?''.
        FRFP reframes it as:
        \begin{itemize}
          \item what explicit semantics does the tool implement?
          \item where are the boundaries between scores and decisions?
          \item how are safe horizons and audits structured?
          \item which institutional operators are responsible for aligning
                explicit tools with tacit norms of justice?
        \end{itemize}
\end{enumerate}

In this sense, COMPAS is less an anomaly than a paradigm case: it exhibits
the characteristic tensions that FRFP is designed to surface and constrain
whenever explicit AI systems are embedded in high-stakes governance
processes.
